% !PRINTED none
% Abschnitt-title page (part divider): no running head/page number visible in
% the scan for this page; §. 169 opens on it.

\begin{center}
\gesperrt{Einundzwanzigster Abschnitt.}

\gesperrt{Kreistheilungsk\"orper.}
\end{center}

\begin{center}
\rule{4cm}{0.4pt}
\end{center}

\begin{center}
\S.\,169.

\gesperrt{Zerlegung der Primzahl $q$ in Primfactoren im Kreistheilungsk\"orper $\Omega_{q^{\varkappa}}$.}
\end{center}

Wir machen eine zweite Anwendung der allgemeinen Theorie der algebraischen Zahlen auf die Kreistheilungsk\"orper, und gewinnen dadurch das Mittel, die Theorie der \gesperrt{Abel'schen} Zahlk\"orper \"uberhaupt zu einem sch\"onen Abschlusse zu bringen.

Die Betrachtungen, die wir zun\"achst anzustellen haben, lassen sich mit kleinen Modificationen auf jeden vollen Kreistheilungsk\"orper (\S.\,18) anwenden. Der Einfachheit halber beschr\"anken wir uns hier aber auf den Fall, dass der Grad der Einheitswurzel eine Primzahlpotenz ist, der f\"ur die beabsichtigte Anwendung ausreicht.

Es sei $q$ eine nat\"urliche Primzahl (mit Einschluss von 2) und
\weq{1}{m = q^{\varkappa}}
eine Potenz von $q$, deren positiver Exponent $\varkappa$ f\"ur $q=2$ gr\"osser als 1 vorausgesetzt wird.

Es sei ferner $r$ eine primitive $m^{\text{te}}$ Einheitswurzel und $\Omega_m$ der volle Kreistheilungsk\"orper f\"ur den Exponenten $m$, dessen Grad
\weq{2}{\mu = \varphi(m) = q^{\varkappa-1}(q-1)}
ist (Bd. I, \S.\,134).

Wir bezeichnen durchweg mit $n$ die $\mu$ Zahlen eines vollen Restsystems f\"ur den Modul $m$ mit Ausschluss der durch $q$ theilbaren Zahlen, und es sind dann die $\mu$ Zahlen $r^n$ die Wurzeln der irreduciblen Gleichung $\mu^{\text{ten}}$ Grades $f(x)=0$, worin
\weq{3}{f(x) = \dfrac{x^m-1}{x^{\frac{m}{q}}-1} = x^{q^{\varkappa-1}(q-1)} + x^{q^{\varkappa-1}(q-2)} + \cdots + 1}
%CONT>>
